\section{模型、对称性与拓扑分类}
\label{sec:model}

\subsection{二聚化 XXZ 链}

考虑具有开边界的 $L$ 格点自旋 $1/2$ 链，以 $X_j,Y_j,Z_j$ 表示 Pauli 算符，晶格常数取一：
\begin{equation}
H(s,\delta)=\sum_{j=1}^{L-1}t_j
\left[e^{-\delta}(X_jX_{j+1}+Y_jY_{j+1})
+e^\delta Z_jZ_{j+1}\right],
\label{eq:H}
\end{equation}
其中 $t_{2r-1}=1-s$、$t_{2r}=s$，且 $0\leq s\leq1$。
交替键沿用 SSH 模型的二聚化结构~\cite{su1979solitons}。
自由费米子能带分类~\cite{ryu2010topological} 提供非相互作用参照；本文的自旋自由度和纵向相互作用对应多体对称性分类。

与文献~\cite{elben2020manybody} 的扩展玻色 SSH 模型对应时，其奇、偶键系数分别为 $J=2(1-s)e^{-\delta}$ 与 $J'=2se^{-\delta}$，各向异性为
\begin{equation}
\Delta=\frac{J_z}{J_{xx}}=e^{2\delta}.
\label{eq:anisotropy}
\end{equation}
因此，$\delta=0$ 为各向同性点，$\delta<0$ 为易面区，$\delta>0$ 为易轴区。这一对应也将本文模型与里德伯原子实现的相互作用玻色 SSH 链联系起来~\cite{deleseleuc2019observation}。除非另有说明，取 $L=8$。

\subsection{对称性与变分扇区}

哈密顿量与总磁化 $Z_{\rm tot}=\sum_jZ_j$、全局自旋翻转 $G_x=\prod_jX_j$ 以及键中心反射 $R:j\mapsto L+1-j$ 对易，也满足自旋时间反演对称性。
$\delta\neq0$ 时连续自旋旋转对称性为 $U(1)$，$\delta=0$ 时扩展为 $SU(2)$。
本文所用链长均为四的倍数，三个参考初态均满足 $Z_{\rm tot}=0$、$G_x=+1$、$R=+1$。
守恒算符
\begin{equation}
P=Z_{\rm tot}^2-G_x-R
\end{equation}
将三个条件合并：$P=-2$ 等价于同时满足这些条件。理想变分电路保持该扇区。

\subsection{两种二聚化为何属于不同相}

对称保护拓扑（SPT）相以有能隙态在保持对称性的局域变换下是否等价来定义~\cite{gu2009tensor,chen2010local}。
对本文的现场对称性，自然的晶胞为 $(2r-1,2r)$，晶胞总自旋为 $T_r^a=(A_{2r-1}+A_{2r})/2$，其中 $a=x,y,z$ 分别对应 $A=X,Y,Z$。
晶胞空间同时包含自旋零和自旋一扇区。
晶胞旋转 $u_x=e^{i\pi T_r^x}$ 与 $u_z=e^{i\pi T_r^z}$ 相互对易，生成 $D_2\simeq\mathbb Z_2\times\mathbb Z_2$。

对于对称且可注入的矩阵乘积态（MPS），对称性在虚拟键上的作用满足
\begin{equation}
\sum_n(u_g)_{mn}A^n=e^{i\phi_g}V_g^\dagger A^mV_g,
\qquad V_gV_h=\omega(g,h)V_{gh}.
\label{eq:mps-symmetry}
\end{equation}
对幺正现场对称群 $G$，不同投影表示由 $H^2(G,U(1))$ 标记，并区分对称有能隙相~\cite{chen2011classification,schuch2011classifying}。
MPS 与 Schmidt 分解框架见综述~\cite{schollwock2011density,cirac2021matrix}。
对 $D_2$，规范不变量关系
\begin{equation}
V_xV_z=\eta V_zV_x,\qquad \eta=\pm1
\label{eq:projective-index}
\end{equation}
区分平庸类与 Haldane 型非平庸类~\cite{pollmann2012detection}。

$s=0$ 时单态位于晶胞内部，构成晶胞直积态，$\eta=+1$；$s=1$ 时体单态跨越晶胞边界，留下具有 $\eta=-1$ 投影表示的边缘自旋 $1/2$。
本文的拓扑参考态将两端自旋组成一个单态，从开链边缘态流形中选取一个态。
这一二聚化极限展示了对称性分数化~\cite{verresen2017onedimensional}。
键反演或时间反演也能独立保护 Haldane 型区分~\cite{pollmann2012symmetry}。
纠缠谱概念最初见文献~\cite{li2008entanglement}，其在一维情形中的简并受到上述对称性保护~\cite{pollmann2010entanglement}。
这些分类结论适用于保护对称性未自发破缺的有能隙体相。AFM 区则在热力学极限自发破缺自旋翻转对称性；有限链的对称猫态仍可保持该对称性，同时具有很强的交错关联。

\subsection{部分反射诊断量}

采用文献~\cite{elben2020manybody} 的玻色部分反射诊断量。
局部空间变换也用于构造费米 SPT 多体不变量~\cite{shiozaki2017manybody}。
对中央区间 $I=I_1\cup I_2$，以 $\rho_I$ 表示约化密度矩阵，$\mathcal R_I$ 将区间内格点逆序，则
\begin{equation}
\tilde Z_{\mathcal R}
=\frac{\operatorname{Tr}(\rho_I\mathcal R_I)}
{\sqrt{\left[\operatorname{Tr}(\rho_{I_1}^2)
+\operatorname{Tr}(\rho_{I_2}^2)\right]/2}}.
\label{eq:reflection-marker}
\end{equation}
$L=8$ 时取 $I=(3,4,5,6)$，左右两半分别为 $(3,4)$ 和 $(5,6)$。
在对称有能隙区，趋于量子化取值 $+1$、$-1$ 要求区间充分大于关联长度；相应模型的 AFM 极限为零~\cite{elben2020manybody}。
本文使用四格点区间计算有限尺寸相图。

式~\eqref{eq:reflection-marker} 由数值波函数计算。
该表达式结合了 $\mathcal R_I$ 的非对角矩阵元与子系统纯度。
配合适当基旋转的随机测量提供了获取这些量的实验途径~\cite{elben2022toolbox}。
实验扫描采用第~\ref{sec:heatmaps} 节构造的计算基关联诊断量 $Q$。
